\section{Taxonomy of Repository Representation Paradigms}\label{sec:taxonomy}

We propose a seven-dimension taxonomy characterising how LLM-based systems
represent code repositories.  The central dimension is the \emph{Paradigm}
(PAR), which admits seven classes:

\begin{description}
  \item[\textbf{PAR-1}] \emph{Dense Embedding Retrieval} --- repository chunks
    encoded as dense vectors; similarity search drives context selection.
  \item[\textbf{PAR-2}] \emph{Cross-File RAG} --- retrieval-augmented generation
    over multi-file repository contents.
  \item[\textbf{PAR-3}] \emph{Graph-Based Representation} --- structural
    relationships (AST, CFG, call graph, DFG, knowledge graph) form the primary
    context substrate.
  \item[\textbf{PAR-4}] \emph{Execution-Based / Dynamic Representation} ---
    runtime traces, test execution results, and coverage data constitute the
    context.
  \item[\textbf{PAR-5}] \emph{Memory-Augmented Representation} --- persistent
    cross-session memory \emph{is} the repository context.
  \item[\textbf{PAR-6}] \emph{Compression \& Summarisation} --- the repository
    is represented through condensed natural-language or structural summaries.
  \item[\textbf{PAR-7}] \emph{Agentic Exploration} --- autonomous agents
    \emph{actively} navigate the repository, constructing context on demand.
\end{description}
